% chapters/requirements/rf_instances.tex
% RF-7: Instancias de examen
% ============================================================================

\item \textbf{Instancias de examen.} \label{rf:instances}
Una instancia de examen es la materialización de un examen para un alumno concreto: el
\acs{pdf} escaneado, sus metadatos (alumno, modelo, emisor, fecha de recepción), sus
anotaciones y sus calificaciones. Cada instancia sigue un ciclo de vida con máquina de
estados que regula las transiciones y las operaciones permitidas en cada estado. Las
incidencias derivadas de la ingesta y el \acs{ocr} se modelan como un \textit{flag}
booleano (\texttt{has\_issues}) ortogonal al estado, de forma que una instancia puede
avanzar en su ciclo de vida mientras tenga incidencias pendientes, salvo para la
transición a publicado, que las requiere resueltas. Los estados del ciclo de vida son:
\texttt{RECEIVED} (recibido), \texttt{QUEUED} (en cola), \texttt{PENDING\_GRADING}
(pendiente de corrección), \texttt{GRADED} (corregido), \texttt{PUBLISHED} (publicado),
\texttt{IN\_REVIEW} (en revisión), \texttt{PENDING\_REVIEW} (pendiente de revisión),
\texttt{FINALIZED} (finalizado) y \texttt{ARCHIVED} (archivado).

\begin{functional}

    % RF-7.1 ----------------------------------------------------------------
    \item \textbf{Creación automática de instancia.} \label{rf:instances:create}
    El sistema debe crear automáticamente una instancia tras la ingesta y el procesamiento
    \acs{ocr} de un \acs{pdf}.

    \textbf{Entrada:} evento interno disparado por la finalización del procesamiento
    \acs{ocr} de un \acs{pdf} ingestado.

    \textbf{Proceso:} el sistema crea un registro de instancia con los datos extraídos por
    el \acs{ocr} o de la petición (alumno identificado, modelo detectado, etc.), la
    referencia al \acs{pdf} almacenado en MinIO, la fecha de recepción y el emisor. El
    estado inicial es \texttt{RECEIVED}. Si algún atributo obligatorio (asociado a una zona
    de reconocimiento) no pudo extraerse, se activa el \textit{flag}
    \texttt{has\_issues=True}. Si el alumno identificado ya tiene una instancia para ese
    examen, la nueva queda sin alumno asignado y con \texttt{has\_issues=True}.

    \textbf{Salida:} instancia creada en la base de datos, con o sin incidencias.

    % RF-7.2 ----------------------------------------------------------------
    \item \textbf{Modificación de instancia.} \label{rf:instances:update}
    El sistema debe permitir al corrector asignado o al coordinador modificar los atributos
    de una instancia.

    \textbf{Entrada:} petición PATCH al recurso de la instancia con los campos a modificar
    (alumno asociado, modelo, emisor y otros atributos).

    \textbf{Proceso:} el sistema valida los datos y aplica los cambios. Si se asigna o
    cambia el alumno, se verifica que no exista ya otra instancia de ese alumno para el
    mismo examen. Si la modificación resuelve todas las incidencias pendientes (por ejemplo,
    se asigna un alumno a una instancia sin identificar y se completan todos los atributos
    obligatorios), el sistema desactiva automáticamente el \textit{flag}
    \texttt{has\_issues}. La asignación de modelo también puede realizarse manualmente
    mediante este recurso, además de por \acs{ocr}.

    \textbf{Salida:} respuesta con los datos actualizados de la instancia.

    % RF-7.3 ----------------------------------------------------------------
    \item \textbf{Consulta de instancia.} \label{rf:instances:detail}
    El sistema debe permitir consultar la información de una instancia, incluyendo el acceso
    al \acs{pdf}, anotaciones y calificaciones.

    \textbf{Entrada:} petición GET al recurso de la instancia.

    \textbf{Proceso:} el sistema recupera los datos de la instancia según el rol y los
    permisos del usuario solicitante. El corrector solo accede al \acs{pdf} de las
    instancias que tiene asignadas (vía \texttt{AssignmentRule}). El estudiante solo accede
    a su propia instancia y solo durante los períodos de publicación o revisión, limitado a
    los permisos configurados en el examen (\cref{rf:exams:student_perms}).

    \textbf{Salida:} respuesta con los datos de la instancia adaptados al nivel de acceso
    del usuario.

    % RF-7.4 ----------------------------------------------------------------
    \item \textbf{Listado de instancias con filtros.} \label{rf:instances:list}
    El sistema debe permitir listar las instancias de un examen con filtros avanzados.

    \textbf{Entrada:} petición GET al recurso de instancias del examen con parámetros
    opcionales de filtrado: grupo, modelo, estado, corrector asignado, presencia de
    incidencias (\texttt{has\_issues}), y paginación.

    \textbf{Proceso:} el sistema aplica los filtros indicados y devuelve los resultados
    paginados. Un corrector solo ve las instancias que tiene asignadas.

    \textbf{Salida:} respuesta con la lista paginada de instancias y metadatos de
    paginación.

    % RF-7.5 ----------------------------------------------------------------
    \item \textbf{Máquina de estados del ciclo de vida.} \label{rf:instances:fsm}
    El sistema debe gestionar las transiciones de estado de cada instancia según la máquina
    de estados definida.

    \textbf{Entrada:} petición PATCH al recurso de la instancia solicitando una transición
    de estado, o evento interno que la desencadena.

    \textbf{Proceso:} el sistema valida que la transición solicitada es legal según la
    máquina de estados. Las transiciones válidas hacia delante son:
    \texttt{RECEIVED}~$\rightarrow$~\texttt{QUEUED}~$\rightarrow$~%
    \texttt{PENDING\_GRADING}~$\rightarrow$~\texttt{GRADED}~$\rightarrow$~%
    \texttt{PUBLISHED}~$\rightarrow$~\texttt{IN\_REVIEW}~$\rightarrow$~%
    \texttt{FINALIZED}~$\rightarrow$~\texttt{ARCHIVED}.
    Desde \texttt{IN\_REVIEW}, las instancias con solicitud de revisión transicionan a
    \texttt{PENDING\_REVIEW} en lugar de a \texttt{FINALIZED} (véase
    \cref{rf:review:close}). Los retrocesos permitidos son:
    \texttt{GRADED}~$\rightarrow$~\texttt{PENDING\_GRADING},
    \texttt{PUBLISHED}~$\rightarrow$~\texttt{GRADED} y
    \texttt{FINALIZED}~$\rightarrow$~\texttt{PENDING\_REVIEW}.
    La transición a \texttt{PENDING\_GRADING}~$\rightarrow$~\texttt{GRADED} requiere
    acción explícita del corrector o el coordinador; no se ejecuta automáticamente aunque
    todos los problemas tengan calificación. La transición a \texttt{PUBLISHED} requiere
    que la instancia no tenga incidencias pendientes (\texttt{has\_issues=False}). La
    transición a \texttt{ARCHIVED} solo se produce por cambio de curso (véase
    \cref{rf:courses:transition}). Se registra cada transición y su fecha en el
    \textit{log} de auditoría. Cualquier transición ilegal se rechaza.

    \textbf{Salida:} respuesta con el nuevo estado de la instancia, o error \acs{http}~409
    si la transición no es válida.

    % RF-7.6 ----------------------------------------------------------------
    \item \textbf{Detección automática de incidencias.} \label{rf:instances:issues}
    El sistema debe detectar automáticamente incidencias derivadas de la ingesta y el
    procesamiento \acs{ocr}.

    \textbf{Entrada:} evento interno disparado tras la creación de una instancia o la
    finalización del procesamiento \acs{ocr}.

    \textbf{Proceso:} el sistema evalúa los siguientes tipos de incidencias: instancia sin
    alumno identificado, dos instancias asignadas al mismo alumno (la segunda queda sin
    asignar) y atributos obligatorios (definidos mediante zonas de reconocimiento) no
    extraídos. Las instancias afectadas se marcan con \texttt{has\_issues=True}. Una vez
    procesados todos los \acs{pdf} del examen, el sistema puede además detectar alumnos
    convocados sin instancia.

    \textbf{Salida:} instancias marcadas con incidencia según corresponda.

    % RF-7.7 ----------------------------------------------------------------
    \item \textbf{Resolución manual de incidencias.} \label{rf:instances:resolve}
    El sistema debe permitir al coordinador o al corrector resolver incidencias manualmente.

    \textbf{Entrada:} petición PATCH al recurso de la instancia afectada, proporcionando los
    datos que resuelven la incidencia (por ejemplo, asignar un alumno a una instancia sin
    identificar, corregir el modelo asignado).

    \textbf{Proceso:} el sistema aplica la corrección y evalúa si todas las incidencias de
    la instancia quedan resueltas. Si todos los atributos obligatorios están completos y no
    hay conflictos, el sistema desactiva el \textit{flag} \texttt{has\_issues}.

    \textbf{Salida:} respuesta con los datos actualizados de la instancia y el estado del
    \textit{flag} de incidencias.

    % RF-7.8 ----------------------------------------------------------------
    \item \textbf{Verificación de cobertura de corrección.}
    \label{rf:instances:coverage}
    El sistema debe verificar que todas las instancias de un examen tienen al menos un
    corrector asignado antes de permitir la publicación.

    \textbf{Entrada:} petición GET al recurso de verificación de cobertura del examen, o
    comprobación automática al solicitar la transición a estado \texttt{PUBLISHED}.

    \textbf{Proceso:} el sistema evalúa cada instancia y cada problema para determinar si
    tienen al menos un corrector asignado mediante \texttt{AssignmentRule}. Identifica las
    instancias y problemas sin cobertura.

    \textbf{Salida:} respuesta con el resultado de la verificación: cobertura completa o
    lista de instancias y problemas sin corrector.

    % RF-7.9 ----------------------------------------------------------------
    \item \textbf{Publicación masiva de notas.} \label{rf:instances:publish}
    El sistema debe permitir al coordinador publicar las calificaciones de un examen,
    cambiando el estado de todas las instancias corregidas.

    \textbf{Entrada:} petición POST al recurso de publicación del examen con la opción de
    notificar a los alumnos.

    \textbf{Proceso:} el sistema verifica que todas las instancias están en estado
    \texttt{GRADED}, que no hay incidencias sin resolver (\texttt{has\_issues=False} en
    todas) y que la cobertura de corrección es completa. Si se cumplen las condiciones,
    transiciona todas las instancias al estado \texttt{PUBLISHED}. Si se solicita
    notificación, se encolan las notificaciones de publicación de notas a los alumnos. Si
    hay instancias sin corregir o con incidencias, la publicación se rechaza indicando el
    motivo.

    \textbf{Salida:} respuesta de confirmación con el número de instancias publicadas, o
    error con el detalle de los impedimentos.

    % RF-7.10 ---------------------------------------------------------------
    \item \textbf{Eliminación de instancia.} \label{rf:instances:delete}
    El sistema debe permitir al coordinador o al gestor eliminar una instancia de examen.

    \textbf{Entrada:} petición DELETE al recurso de la instancia.

    \textbf{Proceso:} el sistema elimina la instancia, sus anotaciones, sus calificaciones
    y el fichero \acs{pdf} asociado en MinIO. Se registra la eliminación en el \textit{log}
    de auditoría.

    \textbf{Salida:} respuesta con código \acs{http}~204.

    % RF-7.11 ---------------------------------------------------------------
    \item \textbf{Descarga de instancia.} \label{rf:instances:download}
    El sistema debe permitir al coordinador o al gestor descargar el \acs{pdf} de una
    instancia con y sin anotaciones.

    \textbf{Entrada:} petición GET al recurso de descarga de la instancia.

    \textbf{Proceso:} el sistema genera una \acs{url} temporal con expiración hacia el
    \acs{pdf} almacenado en MinIO (véase \cref{rf:sec:presigned}). Se registra el acceso
    al \acs{pdf} en el \textit{log} de auditoría.

    \textbf{Salida:} respuesta con la \acs{url} temporal para la descarga del fichero.

    % RF-7.12 ---------------------------------------------------------------
    \item \textbf{Retroceso de estado.} \label{rf:instances:rollback}
    El sistema debe permitir al coordinador retroceder el estado de una instancia para
    corregir errores.

    \textbf{Entrada:} petición PATCH al recurso de la instancia solicitando una transición
    de retroceso.

    \textbf{Proceso:} el sistema valida que la transición de retroceso es una de las
    permitidas (\texttt{GRADED}~$\rightarrow$~\texttt{PENDING\_GRADING},
    \texttt{PUBLISHED}~$\rightarrow$~\texttt{GRADED},
    \texttt{FINALIZED}~$\rightarrow$~\texttt{PENDING\_REVIEW}) y la ejecuta. Se registra
    el retroceso en el \textit{log} de auditoría.

    \textbf{Salida:} respuesta con el nuevo estado de la instancia.

    % RF-7.13 ---------------------------------------------------------------
    \item \textbf{Almacenamiento de \acs{pdf} de instancias en MinIO.}
    \label{rf:instances:storage}
    El sistema debe almacenar los ficheros \acs{pdf} escaneados de las instancias en el
    servicio de almacenamiento de objetos.

    \textbf{Entrada:} evento interno disparado por la ingesta de un \acs{pdf}.

    \textbf{Proceso:} el fichero se almacena en MinIO en un \dfn{bucket} dedicado a
    instancias, organizado por organización, asignatura y examen. Se almacenan los
    metadatos del fichero (tamaño, número de páginas, fecha de ingesta) en la base de
    datos.

    \textbf{Salida:} confirmación del almacenamiento con la referencia interna al objeto en
    MinIO.

\end{functional}
